\begin{remark}[Set theoretic issues]
\label{remark-set-theoretic-LC}
The category $\textit{LC}$ is a ``big'' category as its objects form
a proper class. Similarly, the coverings form a proper class.
Let us define the {\it size} of a topological space $X$ to be the
cardinality of the set of points of $X$. Choose a function
$Bound$ on cardinals, for example as in
Sets, Equation (\ref{sets-equation-bound}).
Finally, let $S_0$ be an initial set of objects of $\textit{LC}$,
for example $S_0 = \{(\mathbf{R}, \text{euclidean topology})\}$.
Exactly as in Sets, Lemma \ref{sets-lemma-construct-category}
we can choose a limit ordinal $\alpha$ such that
$\textit{LC}_\alpha = \textit{LC} \cap V_\alpha$
contains $S_0$ and is preserved under all countable limits and
colimits which exist in $\textit{LC}$. Moreover, if $X \in \textit{LC}_\alpha$
and if $Y \in \textit{LC}$ and
$\text{size}(Y) \leq Bound(\text{size}(X))$, then $Y$ is isomorphic
to an object of $\textit{LC}_\alpha$.
Next, we apply Sets, Lemma \ref{sets-lemma-coverings-site}
to choose set $\text{Cov}$ of qc covering on $\textit{LC}_\alpha$
such that every qc covering in $\textit{LC}_\alpha$ is
combinatorially equivalent to a covering this set.
In this way we obtain a site $(\textit{LC}_\alpha, \text{Cov})$
which we will denote $\textit{LC}_{qc}$.
\end{remark}